\documentclass[11pt,a4paper]{article}
\usepackage[T1]{fontenc}
\usepackage{lmodern}
\usepackage{amsmath,amssymb,amsthm,mathtools}
\usepackage[margin=27mm]{geometry}
\usepackage{microtype}
\usepackage[hidelinks]{hyperref}
\hypersetup{pdftitle={Residual conductivity at isotropic Dirac points},
  pdfsubject={A proof for periodic Schrodinger operators},
  pdfkeywords={conductivity, Dirac points, linear response, periodic operators}}
\numberwithin{equation}{section}

\newcommand{\R}{\mathbb R}
\newcommand{\C}{\mathbb C}
\newcommand{\Z}{\mathbb Z}
\newcommand{\T}{\mathbb T}
\newcommand{\one}{\mathbf 1}
\newcommand{\cR}{\mathcal R}
\newcommand{\cB}{\mathcal B}
\newcommand{\cD}{\mathcal D}
\newcommand{\cH}{\mathcal H}
\newcommand{\trv}{\underline{\operatorname{Tr}}}
\newcommand{\dd}{\,\mathrm d}
\newcommand{\ii}{\mathrm i}
\DeclareMathOperator{\Tr}{Tr}
\DeclareMathOperator{\Ran}{Ran}
\DeclareMathOperator{\Rea}{Re}
\DeclareMathOperator{\Ima}{Im}
\DeclareMathOperator{\dist}{dist}
\newtheorem{theorem}{Theorem}[section]
\newtheorem{proposition}[theorem]{Proposition}
\newtheorem{lemma}[theorem]{Lemma}
\theoremstyle{remark}
\newtheorem{remark}[theorem]{Remark}

\title{Residual conductivity of a periodic Schr\"odinger operator\\
with isotropic Dirac points}
\author{}
\date{}

\begin{document}
\maketitle

\begin{abstract}
We prove a residual conductivity formula for a periodic Schr\"odinger
operator in two dimensions, describing non-interacting spinless electrons.
Its Fermi surface consists
of finitely many isolated, isotropic, twofold conical crossings. If there
are $M$ such crossings, the Ces\`aro mean of the current response to a
constant electric field, with the linear-response limit taken first, is
$M/16$ times the Euclidean conductivity tensor in atomic units.
The proof derives the response from the Stark evolution, establishes an
integrable band sum rule that cancels the term linear in time, and reduces
each crossing to a two-dimensional matrix calculation. In particular,
the assumptions on the band energies imply the required velocity matrix
elements; these are not additional hypotheses.
\end{abstract}

\section{Setting and theorem}

Let $\cR\subset\R^2$ be a lattice, let $\Omega$ be a fundamental cell, and
let $\cB$ be a fundamental cell of the reciprocal lattice
\[
 \cR^*=\{K\in\R^2:K\cdot R\in2\pi\Z\text{ for all }R\in\cR\}.
\]
We identify quasimomenta modulo $\cR^*$ and write
$\T_*^2=\R^2/\cR^*$. In particular,
\begin{equation}\label{eq:volumes}
 |\Omega|\,|\cB|=(2\pi)^2.
\end{equation}
All spatial scalar products and norms are Euclidean. We use atomic units
and consider one spin species.

For a real potential $V\in H^1_{\mathrm{per}}(\Omega)$, put
\begin{equation}\label{eq:hamiltonians}
 H=\tfrac12p^2+V,\qquad p=-\ii\nabla,
 \qquad H_k=\tfrac12(p+k)^2+V.
\end{equation}
Here $H$ acts on $L^2(\R^2)$ and $H_k$ acts on
$\cH=L^2_{\mathrm{per}}(\Omega)$, with domain
$H^2_{\mathrm{per}}(\Omega)$. Its eigenvalues, counted with multiplicity,
are $\lambda_1(k)\leq\lambda_2(k)\leq\cdots$; set
$\lambda_0(k)=-\infty$.

Fix $\mu_{\mathrm F}\in\R$. The hypotheses are the following.
\begin{enumerate}
\item There are an integer $N\geq1$ and distinct points
 $\cD=\{k_1,\ldots,k_M\}\subset\T_*^2$ such that
 \begin{equation}\label{eq:gap}
  \lambda_N(k)<\mu_{\mathrm F}<\lambda_{N+1}(k)
  \qquad(k\in\T_*^2\setminus\cD).
 \end{equation}
\item At each $k_i$, precisely two bands meet at the Fermi energy:
 \begin{equation}\label{eq:double}
  \lambda_{N-1}(k_i)<\lambda_N(k_i)=\mu_{\mathrm F}
  =\lambda_{N+1}(k_i)<\lambda_{N+2}(k_i).
 \end{equation}
\item For each $i$ there is $v_i>0$ such that, as $q\to0$ in $\R^2$,
 \begin{equation}\label{eq:cones}
 \begin{split}
  \lambda_N(k_i+q)&=\mu_{\mathrm F}-v_i|q|+O(|q|^2),\\
  \lambda_{N+1}(k_i+q)&=\mu_{\mathrm F}+v_i|q|+O(|q|^2).
 \end{split}
 \end{equation}
\end{enumerate}
The error bounds in \eqref{eq:cones} hold for all sufficiently small $q$,
uniformly in its direction.

Let $e_\beta,e_\alpha\in\R^2$ denote the field and current directions;
they need not be unit vectors. For $\varepsilon\geq0$, define
\begin{equation}\label{eq:stark}
 \gamma_0=\one_{(-\infty,\mu_{\mathrm F}]}(H),\qquad
 H_\varepsilon=H-\varepsilon x\cdot e_\beta,\qquad
 J_\alpha=p\cdot e_\alpha,
\end{equation}
where $H_\varepsilon$ is the self-adjoint closure of its restriction to
$C_c^\infty(\R^2)$. This closure exists, as shown below. Set
\begin{equation}\label{eq:current}
 \gamma_\varepsilon(t)=e^{-\ii tH_\varepsilon}\gamma_0
                            e^{\ii tH_\varepsilon},\qquad
 j_\varepsilon(t)=\trv(J_\alpha\gamma_\varepsilon(t)),
 \qquad t\geq0,
\end{equation}
where, for the periodic locally trace-class operators in question,
\begin{equation}\label{eq:trace-definition}
 \trv(A)=\frac1{|\Omega|}\Tr(\one_\Omega A\one_\Omega).
\end{equation}

\begin{theorem}\label{thm:main}
Under \eqref{eq:gap}--\eqref{eq:cones}, the current
\eqref{eq:current} is well defined for every $\varepsilon,t\geq0$,
and $j_0(t)=0$. For every fixed $t\geq0$ the limit
\[
 \sigma_{\alpha\beta}(t)
   =\lim_{\varepsilon\downarrow0}\frac{j_\varepsilon(t)}{\varepsilon}
\]
exists. Its Ces\`aro mean satisfies
\begin{equation}\label{eq:main}
 \lim_{T\to\infty}\frac1T\int_0^T
   \left(\lim_{\varepsilon\downarrow0}
           \frac{j_\varepsilon(t)}{\varepsilon}\right)\dd t
       =\frac{M}{16}\,e_\alpha\cdot e_\beta.
\end{equation}
\end{theorem}

The theorem is conditional on the stated band structure. No lattice
symmetry beyond periodicity is imposed, and no simplicity assumption is
made on bands away from the Fermi crossings. The order of limits in
\eqref{eq:main} is part of the assertion.

\section{Operator framework and the exact current}

We first justify the dynamics, the trace, and the differentiations used
later. Standard operator results are cited from Reed--Simon; the estimates
needed at the conical crossings will be proved explicitly.

\subsection{Domains and the Stark operator}

\begin{lemma}\label{lem:domains}
The operator $H$ is self-adjoint on $H^2(\R^2)$, and each $H_k$ is
self-adjoint on $H^2_{\mathrm{per}}(\Omega)$ with compact resolvent.
For $k$ in a compact set, its graph norm is uniformly equivalent to the
$H^2$ norm. The family $k\mapsto H_k$ is analytic with this common domain.
Moreover, $H-\varepsilon x\cdot e_\beta$ is essentially self-adjoint on
$C_c^\infty(\R^2)$ for every real $\varepsilon$.
\end{lemma}

\begin{proof}
Periodicity and $V\in L^2(\Omega)$ imply that $V$ is uniformly locally
$L^2$. The localization estimate of
\cite[Theorem XIII.96]{RS4} gives, for every $\eta>0$,
\begin{equation}\label{eq:potential-bound}
 \|Vu\|_2\leq\eta\|\Delta u\|_2+C_\eta\|u\|_2.
\end{equation}
The same estimate holds on the torus. Also,
\begin{equation}\label{eq:form-bound}
 |\langle u,Vu\rangle|
       \leq\eta\|\nabla u\|_2^2+C_\eta\|u\|_2^2.
\end{equation}
For completeness, on a bounded cell this follows from H\"older's
inequality and the two-dimensional interpolation estimate
$\|u\|_4^2\leq C\|u\|_2(\|\nabla u\|_2+\|u\|_2)$.
A partition of unity by translates of bounded cells, with bounded
overlap, proves the whole-space version. The Kato--Rellich theorem
\cite[Theorem X.12]{RS2}, applied also to $k\cdot p$, proves the domain
and graph-norm claims. Compactness follows from the compact embedding
$H^2_{\mathrm{per}}\hookrightarrow L^2_{\mathrm{per}}$.
The polynomial dependence on $k$, as a bounded map from $H^2$ to $L^2$,
gives analyticity. Equivalently, a local Neumann series for the resolvent
gives analyticity in the several variables $k_1,k_2$; see also
\cite[Section XII.2 and Section XIII.16]{RS4}.

Here are details for the unbounded electric potential. On
$C_c^\infty(\R^2)$ introduce
\[
 \mathcal N=-\Delta+2V+|x|^2+b.
\]
The harmonic oscillator $-\Delta+|x|^2+1$ has this space as a core
by \cite[Theorem X.28]{RS2}. Its graph norm controls both $\Delta$
and $|x|^2$, as follows by integration by parts from
\[
 \|\Delta u\|_2^2+\||x|^2u\|_2^2
    \leq\|(-\Delta+|x|^2)u\|_2^2+4\|u\|_2^2.
\]
Consequently \eqref{eq:potential-bound} and Kato--Rellich show that
$\mathcal N$ is self-adjoint on the oscillator domain and has the same
core. Choose $b$ large enough that $\mathcal N\geq1$. By
\eqref{eq:form-bound}, its form controls
$\|\nabla u\|_2^2+\|xu\|_2^2+\|u\|_2^2$.
If $A=H-\varepsilon x\cdot e_\beta$, graph-norm equivalence gives
\[
 \|Au\|_2\leq C_\varepsilon\|\mathcal N u\|_2.
\]
Since $\mathcal N=2H+|x|^2+b$, cancellation of the multiplication
operators in the commutator gives, in the quadratic-form sense on the
core,
\[
 [A,\mathcal N]
   =-\ii(p\cdot x+x\cdot p)-2\ii\varepsilon p\cdot e_\beta.
\]
It follows that
\[
 |\langle Au,\mathcal N u\rangle
           -\langle\mathcal N u,Au\rangle|
   \leq C_\varepsilon\langle u,\mathcal N u\rangle.
\]
The commutator theorem \cite[Theorem X.37]{RS2} now proves essential
self-adjointness of $A$ on $C_c^\infty$. This argument does not
differentiate the possibly unbounded potential $V$.
\end{proof}

We use the Bloch transform normalized by
\begin{equation}\label{eq:bloch}
 (\mathcal U f)(k,x)
   =\sum_{R\in\cR}e^{-\ii k\cdot(x+R)}f(x+R),
 \qquad
 \mathcal U:L^2(\R^2)\longrightarrow
       \int_{\cB}^{\oplus}\cH\,\frac{\dd k}{|\cB|}.
\end{equation}
It is unitary and decomposes $H$ into $H_k$ and $J_\alpha$ into
\begin{equation}\label{eq:velocity}
 v_\alpha(k)=(p+k)\cdot e_\alpha=\partial_\alpha H_k,
 \qquad \partial_\alpha=e_\alpha\cdot\nabla_k.
\end{equation}
These are the usual periodic-operator decompositions; see
\cite[Theorems XIII.97--XIII.99]{RS4}. Under a reciprocal translation
$K$, the fibers are conjugated by multiplication by $e^{-\ii K\cdot x}$.
All scalar traces below are therefore functions on $\T_*^2$.

Write
\begin{equation}\label{eq:fermi-projection}
 P(k)=\one_{(-\infty,\mu_{\mathrm F}]}(H_k).
\end{equation}
It has rank $N$ almost everywhere. At a point of $\cD$ its actual rank
is $N+1$, but changing it on this finite set does not change any
whole-space operator or integral. When convenient we choose there a
commuting rank-$N$ projection containing the lower bands. Uniform
graph-norm estimates and the bounded occupied energies give
\begin{equation}\label{eq:occupied-bound}
 \sup_{k\in\cB}\|P(k)\|_{L^2\to H^2}<\infty,
 \qquad
 \sup_{k\in\cB}\sum_{n=1}^N\|v_\alpha(k)u_n(k)\|_2^2<\infty,
\end{equation}
where $u_1(k),\ldots,u_N(k)$ is any orthonormal occupied eigenbasis.
No regularity of this choice of eigenvectors is asserted or needed.

\subsection{Temporal gauge}

Let $U_k^\varepsilon(t,s)$ be the propagator solving
\begin{equation}\label{eq:propagator}
 \ii\partial_t U_k^\varepsilon(t,s)
       =H_{k+\varepsilon t e_\beta}U_k^\varepsilon(t,s),
 \qquad U_k^\varepsilon(s,s)=1.
\end{equation}
The common-domain evolution theorem
\cite[Theorem X.70]{RS2} applies: the domain is $H^2_{\mathrm{per}}$,
and the Hamiltonians and their time derivatives are continuous as maps
$H^2_{\mathrm{per}}\to L^2_{\mathrm{per}}$, with uniformly equivalent
graph norms on compact parameter sets. It gives a unitary propagator
preserving $H^2_{\mathrm{per}}$, with locally uniform bounds in that
norm. The invertibility condition in that theorem is obtained by adding
a constant so that the Hamiltonians are at least $1$ on the parameter
set in question, and then removing the resulting scalar phase.
The identical argument on $\R^2$ gives the propagator
$U^\varepsilon(t,s)$ for
\[
 H^{\mathrm g}_\varepsilon(t)
     =\tfrac12(p+\varepsilon t e_\beta)^2+V,
\]
whose Bloch fibers are \eqref{eq:propagator}.

\begin{lemma}\label{lem:gauge}
For $M_t=e^{\ii\varepsilon t x\cdot e_\beta}$ one has
\begin{equation}\label{eq:gauge}
 e^{-\ii tH_\varepsilon}=M_tU^\varepsilon(t,0).
\end{equation}
The current is
\begin{equation}\label{eq:exact-current}
 j_\varepsilon(t)=\frac1{(2\pi)^2}\int_{\cB}
  \Tr\!\left[
   \bigl(v_\alpha(k)+\varepsilon t(e_\alpha\cdot e_\beta)\bigr)
   U_k^\varepsilon(t,0)P(k)U_k^\varepsilon(t,0)^*
       \right]\dd k.
\end{equation}
In particular, the localized operator defining this current is trace
class for every finite $\varepsilon,t$.
\end{lemma}

\begin{proof}
The covariance
$H^{\mathrm g}_\varepsilon(r+s)
       =M_s^*H^{\mathrm g}_\varepsilon(r)M_s$
and uniqueness of the propagator imply
\[
 U^\varepsilon(t+s,s)=M_s^*U^\varepsilon(t,0)M_s.
\]
Thus $W(t)=M_tU^\varepsilon(t,0)$ is a strongly continuous unitary
group. Its derivative at zero on $C_c^\infty$ is
$-\ii(H-\varepsilon x\cdot e_\beta)$.
By Lemma~\ref{lem:domains} its self-adjoint generator must be
$H_\varepsilon$, proving \eqref{eq:gauge}.

Put $\rho^\varepsilon(t)=U^\varepsilon(t,0)\gamma_0
U^\varepsilon(t,0)^*$. This is periodic, with rank-$N$ fiber projections
\[
 \rho_k^\varepsilon(t)
       =U_k^\varepsilon(t,0)P(k)U_k^\varepsilon(t,0)^*.
\]
The operators $U_k^\varepsilon(t,0)P(k):L^2\to H^2$ are uniformly
bounded on compact $(\varepsilon,t)$ sets, by
\eqref{eq:occupied-bound} and the propagator bounds.
The density $\gamma_\varepsilon(t)=M_t\rho^\varepsilon(t)M_t^*$
is also periodic: conjugation of $M_t$ by a lattice translation changes
it only by a scalar phase.

We recall the trace normalization and explain the trace-class assertion.
For a bounded periodic operator $B$ with square-integrable
Hilbert--Schmidt fibers, Bloch Parseval gives
\begin{equation}\label{eq:localized-HS}
 \|\one_\Omega B\|_{\mathfrak S_2}^2
    =\frac1{|\cB|}\int_{\cB}
       \|B(k)\|_{\mathfrak S_2}^2\dd k.
\end{equation}
One can see this by inserting an orthonormal basis of functions supported
in $\Omega$ into $B^*\one_\Omega$; its Bloch transforms are the same
basis multiplied by $e^{-\ii k\cdot x}$. Applying Parseval and then
summing proves \eqref{eq:localized-HS}.
The analogous identity holds for $B\one_\Omega$.
The product of two such localized Hilbert--Schmidt operators is trace
class, and polarization of the same identity gives the fiber trace
formula for these products.

For a periodic projection $Q$ of uniformly bounded fiber rank with
$\sup_k\|Q(k)\|_{L^2\to H^1}<\infty$, both $\one_\Omega J_\alpha Q$ and
$Q\one_\Omega$ are Hilbert--Schmidt. Consequently
\[
 \one_\Omega J_\alpha Q\one_\Omega
     =(\one_\Omega J_\alpha Q)(Q\one_\Omega)
\]
is trace class. The same argument applies after the multiplication by
$M_t$, which only shifts momentum by $\varepsilon t e_\beta$.
For all operators used here it proves
\begin{equation}\label{eq:trace-formula}
 \trv(A)=\frac1{|\Omega|\,|\cB|}
             \int_{\cB}\Tr A(k)\dd k
          =\frac1{(2\pi)^2}\int_{\cB}\Tr A(k)\dd k.
\end{equation}
Finally, $M_t$ commutes with $\one_\Omega$ and
$M_t^*J_\alpha M_t=J_\alpha+\varepsilon t(e_\alpha\cdot e_\beta)$.
Unitary invariance of the localized trace, followed by
\eqref{eq:trace-formula}, gives \eqref{eq:exact-current}.
\end{proof}

\section{A band sum rule at the Fermi energy}

For $k\notin\cD$, choose a measurable orthonormal eigenbasis $u_n(k)$
of $H_k$; such a choice is available as in
\cite[Theorem XIII.98]{RS4}.
For $n\leq N<m$ define
\begin{equation}\label{eq:matrix-elements}
 \Delta_{mn}(k)=\lambda_m(k)-\lambda_n(k)>0,\qquad
 c_{nm}^{\alpha\beta}(k)
   =\langle u_n,v_\alpha u_m\rangle
      \langle u_m,v_\beta u_n\rangle.
\end{equation}
Inner products are conjugate linear in the first variable. Sums over
complete eigenspaces are independent of the chosen basis, including at
degeneracies within the occupied or unoccupied bands. Cauchy--Schwarz
and \eqref{eq:occupied-bound} imply the useful uniform estimate
\begin{equation}\label{eq:matrix-sum-bound}
 \sum_{n\leq N<m}|c_{nm}^{\alpha\beta}(k)|
 \leq
 \left(\sum_{n\leq N}\|v_\alpha u_n\|_2^2\right)^{1/2}
 \left(\sum_{n\leq N}\|v_\beta u_n\|_2^2\right)^{1/2}
 \leq C_{\alpha\beta}.
\end{equation}
In particular, the infinite sums below are absolutely convergent.

Let
\begin{equation}\label{eq:filled-energy}
 F(k)=\sum_{n=1}^N\lambda_n(k),\qquad
 g(k)=\lambda_{N+1}(k)-\lambda_N(k).
\end{equation}
Continuity of the bands, compactness of $\T_*^2$, and
\eqref{eq:cones} give
\begin{equation}\label{eq:gap-lower}
 g(k)\geq c\dist(k,\cD)\quad\hbox{near }\cD,
 \qquad \inf_{\dist(k,\cD)\geq\delta}g(k)>0\quad(\delta>0).
\end{equation}
Thus \eqref{eq:matrix-sum-bound} implies
\begin{equation}\label{eq:integrable-sum}
 \int_{\cB}\sum_{n\leq N<m}
       \frac{|c_{nm}^{\alpha\beta}(k)|}{\Delta_{mn}(k)}\dd k<\infty.
\end{equation}

\begin{proposition}\label{prop:sumrule}
The function $F$ is smooth off $\cD$, has bounded first derivatives
there, and satisfies
\begin{align}
 \partial_\alpha F
     &=\Tr(v_\alpha P),\label{eq:HF}\\
 \partial_\beta\partial_\alpha F
     &=N(e_\alpha\cdot e_\beta)
       -2\sum_{n\leq N<m}
             \frac{\Rea c_{nm}^{\alpha\beta}}{\Delta_{mn}}.
       \label{eq:sumrule-pointwise}
\end{align}
Both derivatives are integrable on the torus and have integral zero.
In particular,
\begin{equation}\label{eq:sumrule-integrated}
 N|\cB|(e_\alpha\cdot e_\beta)
    =2\int_{\cB}\sum_{n\leq N<m}
          \frac{\Rea c_{nm}^{\alpha\beta}}{\Delta_{mn}}\dd k.
\end{equation}
Also $j_0(t)=0$.
\end{proposition}

\begin{proof}
Off $\cD$, the occupied spectral cluster is isolated. The resolvent
contour formula makes $P(k)$ analytic as a map $L^2\to H^2$;
this is the usual Riesz-projection argument of
\cite[Section XII.2]{RS4}. Since $\partial_\beta P$ is off diagonal
relative to $P\oplus(1-P)$, differentiating $F=\Tr(H_kP)$ proves
\eqref{eq:HF}. Its right-hand side is uniformly bounded by
\eqref{eq:occupied-bound}.

Differentiating $[H_k,P]=0$ and $P^2=P$ gives, for $n\leq N<m$,
\[
 \langle u_m,(\partial_\beta P)u_n\rangle
       =-\frac{\langle u_m,v_\beta u_n\rangle}{\Delta_{mn}}.
\]
The opposite off-diagonal block is its adjoint. Differentiating
\eqref{eq:HF}, using
$\partial_\beta v_\alpha=(e_\alpha\cdot e_\beta)1$, proves
\eqref{eq:sumrule-pointwise}. This proof uses the spectral projection
and remains valid at internal band degeneracies.

The right-hand side of \eqref{eq:sumrule-pointwise} is in $L^1$ by
\eqref{eq:integrable-sum}. Remove disjoint disks of radius $\delta$
around the points of $\cD$ and apply the divergence theorem on the
remaining torus. The boundary integral for
$\partial_\beta\partial_\alpha F$ is bounded by
$C\delta\sup|\partial_\alpha F|$, hence tends to zero. This proves
that the integral of the Hessian entry is zero. The same argument for
$\partial_\alpha F$, using boundedness of the continuous function
$F$, proves that its integral is zero. In particular, the conical
singularities create no missing point-supported term in the integrated
sum rule.

At $\varepsilon=0$, $U_k^0(t,0)$ commutes with $P(k)$, so
\eqref{eq:exact-current} and \eqref{eq:HF} give
$j_0(t)=(2\pi)^{-2}\int_{\cB}\partial_\alpha F\dd k=0$.
\end{proof}

\section{Linear response at fixed time}

\begin{proposition}\label{prop:response}
For every fixed $t\geq0$ the linear response exists and is
\begin{equation}\label{eq:response}
 \sigma_{\alpha\beta}(t)
   =\frac{2}{(2\pi)^2}\int_{\cB}
       \sum_{n\leq N<m}\Rea c_{nm}^{\alpha\beta}(k)
       \frac{\sin(t\Delta_{mn}(k))}{\Delta_{mn}(k)^2}\dd k.
\end{equation}
This formula is absolutely integrable for each finite $t$.
\end{proposition}

\begin{proof}
We first justify differentiation with the unbounded current observable.
Fix a bounded time interval and take $|\varepsilon|\leq1$. For a unit
vector $u$ in the occupied subspace, set
\[
 \psi_\varepsilon(t)=U_k^\varepsilon(t,0)u,\qquad
 \psi_0(t)=e^{-\ii tH_k}u.
\]
Since
\begin{equation}\label{eq:perturbation}
 H_{k+\varepsilon s e_\beta}-H_k
       =\varepsilon s v_\beta(k)
          +\tfrac12\varepsilon^2s^2|e_\beta|^2,
\end{equation}
Duhamel's formula and the uniform $H^2$ bounds give
\begin{equation}\label{eq:duhamel-bound}
 \|\psi_\varepsilon-\psi_0\|_2\leq C|\varepsilon|,
 \qquad
 \|\psi_\varepsilon\|_{H^2}+\|\psi_0\|_{H^2}\leq C,
\end{equation}
uniformly in $k$ and $u$. Moreover,
\begin{equation}\label{eq:wave-derivative}
 \frac{\psi_\varepsilon(t)-\psi_0(t)}{\varepsilon}
   \longrightarrow
 \dot\psi(t)=-\ii\int_0^t
       e^{-\ii(t-s)H_k}s v_\beta(k)e^{-\ii sH_k}u\dd s
 \quad\hbox{in }L^2,
\end{equation}
again uniformly for these parameters. To verify the uniformity without
any regular choice of $u(k)$, Duhamel gives
$\|U_k^\varepsilon(t,s)-U_k^0(t,s)\|_{H^2\to L^2}
\leq C|\varepsilon|$.
Interpolation with the $L^2$ bound gives an $O(|\varepsilon|^{1/2})$
bound from $H^1$ to $L^2$. The vector
$v_\beta e^{-\ii sH_k}u$ is uniformly bounded in $H^1$, which proves
\eqref{eq:wave-derivative} by the Duhamel formula with
$U_k^\varepsilon(t,s)$ on the left of the perturbation.

Writing $w=\psi_\varepsilon-\psi_0$, interpolation in
\eqref{eq:duhamel-bound} gives $\|w\|_{H^1}=O(|\varepsilon|^{1/2})$.
The exact expansion
\[
 \langle\psi_\varepsilon,v_\alpha\psi_\varepsilon\rangle
 -\langle\psi_0,v_\alpha\psi_0\rangle
 =2\Rea\langle w,v_\alpha\psi_0\rangle
      +\langle w,v_\alpha w\rangle
\]
has last term $O(|\varepsilon|^{3/2})$. It follows that the current
expectations are differentiable at zero, uniformly in $k$. This also
justifies differentiation under the integral in
\eqref{eq:exact-current}.

In density-matrix notation the derivative is
\begin{equation}\label{eq:density-derivative}
 \dot\rho_k(t)=-\ii\int_0^t s\,
       e^{-\ii(t-s)H_k}[v_\beta(k),P(k)]
       e^{\ii(t-s)H_k}\dd s.
\end{equation}
This expression may equally be read as the finite-rank sum obtained
from \eqref{eq:wave-derivative}. The free propagator preserves $H^1$,
the form domain of $H_k$ shifted by a constant, so the velocity traces
in this formula are well defined.

For $\tau=t-s$, insertion of an eigenbasis in
\eqref{eq:density-derivative} gives
\[
 -\ii\Tr\bigl(v_\alpha e^{-\ii\tau H_k}
             [v_\beta,P]e^{\ii\tau H_k}\bigr)
 =2\sum_{n\leq N<m}
       \Ima\bigl(c_{nm}^{\alpha\beta}e^{-\ii\tau\Delta_{mn}}\bigr).
\]
Since
\[
 \int_0^t s\sin((t-s)\Delta)\dd s
     =\frac{t}{\Delta}-\frac{\sin(t\Delta)}{\Delta^2},\qquad
 \int_0^t s\cos((t-s)\Delta)\dd s
     =\frac{1-\cos(t\Delta)}{\Delta^2},
\]
the derivative of \eqref{eq:exact-current} is
\begin{equation}\label{eq:full-response}
\begin{split}
 \sigma_{\alpha\beta}(t)
  =\frac1{(2\pi)^2}\int_{\cB}\Bigg\{
   &tN(e_\alpha\cdot e_\beta)\\
   &-2\sum_{n\leq N<m}\Rea c_{nm}^{\alpha\beta}
          \left(\frac{t}{\Delta_{mn}}
                    -\frac{\sin(t\Delta_{mn})}{\Delta_{mn}^2}\right)\\
   &+2\sum_{n\leq N<m}\Ima c_{nm}^{\alpha\beta}
             \frac{1-\cos(t\Delta_{mn})}{\Delta_{mn}^2}
             \Bigg\}\dd k.
\end{split}
\end{equation}
All terms are integrable by \eqref{eq:matrix-sum-bound},
\eqref{eq:integrable-sum}, and the bounds
$|\sin(t\Delta)|\leq t\Delta$ and
$1-\cos(t\Delta)\leq t^2\Delta^2/2$.

Realness of $V$ now removes the last line of \eqref{eq:full-response}.
Indeed, complex conjugation $\mathcal C$ satisfies
\begin{equation}\label{eq:time-reversal}
 \mathcal C H_k\mathcal C=H_{-k},\qquad
 \mathcal C v_\alpha(k)\mathcal C=-v_\alpha(-k).
\end{equation}
Eigenvectors can be paired locally at $k$ and $-k$ by $\mathcal C$.
In these paired bases the gaps agree and
$c_{nm}^{\alpha\beta}(-k)=\overline{c_{nm}^{\alpha\beta}(k)}$.
For degenerate eigenvalues this statement is used after summing over
the corresponding eigenspaces, so it is basis independent. The last
line is therefore an odd integrable function on the torus and has
integral zero.

Finally \eqref{eq:sumrule-integrated} cancels exactly the coefficient
of $t$ in \eqref{eq:full-response}. This proves \eqref{eq:response}.
Its asserted integrability follows directly from
$|\sin(t\Delta)|/\Delta^2\leq t/\Delta$ and
\eqref{eq:integrable-sum}.
\end{proof}

\section{The matrix structure forced by an isotropic cone}

Only one transition near each $k_i$ can have a small gap. We now
determine its leading velocity matrix element from the band hypotheses.
Write $q=r\omega$, $r>0$, $\omega\in S^1$, and set
\begin{equation}\label{eq:local-definitions}
 \Delta_i(q)=\lambda_{N+1}(k_i+q)-\lambda_N(k_i+q),\qquad
 G_i^{\alpha\beta}(q)=\Rea c_{N,N+1}^{\alpha\beta}(k_i+q).
\end{equation}

\begin{lemma}\label{lem:cone}
Uniformly in $\omega$ as $r\downarrow0$,
\begin{align}
 \Delta_i(r\omega)&=2v_i r+O(r^2),\label{eq:local-gap}\\
 G_i^{\alpha\beta}(r\omega)
   &=v_i^2\bigl(e_\alpha\cdot e_\beta
            -(e_\alpha\cdot\omega)(e_\beta\cdot\omega)\bigr)+O(r).
       \label{eq:local-matrix}
\end{align}
\end{lemma}

\begin{proof}
Fix a crossing and temporarily suppress $i$. By \eqref{eq:double},
a fixed resolvent contour around $\mu_{\mathrm F}$ defines an analytic
rank-two spectral projection $\Pi(q)$ for small $q$, including $q=0$.
It is analytic as a map $L^2\to H^2$. Choose an isometry
$W_0:\C^2\to\Ran\Pi(0)$ and define
\begin{equation}\label{eq:frame}
 W(q)=\Pi(q)W_0\bigl(W_0^*\Pi(q)W_0\bigr)^{-1/2}.
\end{equation}
For small real $q$ the matrix inside the square root is positive and
close to the identity. Thus $W(q)$ is a real-analytic isometry into
$H^2$, $W(q)W(q)^*=\Pi(q)$, and $W(0)=W_0$.
This explicit frame is a local version of analytic spectral-subspace
reduction; compare \cite[Theorem XII.12]{RS4}. It does not assert
analyticity of the two individual bands at the crossing.

The Hermitian matrix
\[
 h(q)=W(q)^*H_{k_i+q}W(q)
\]
has precisely the two crossing eigenvalues. With respect to fixed
orthonormal coordinates in $\R^2$, Taylor expansion gives
\begin{equation}\label{eq:linear-matrix}
 h(q)=\mu_{\mathrm F}I+q_1A_1+q_2A_2+O(|q|^2),\qquad
 A_j=W_0^*(p_j+(k_i)_j)W_0.
\end{equation}
To obtain the formula for $A_j$, differentiate $h$ at zero and use
$H_{k_i}W_0=\mu_{\mathrm F}W_0$ and
$\partial_j(W^*W)=0$; the two frame-derivative terms cancel.

For any $\omega\in S^1$, the eigenvalues of
$\omega_1A_1+\omega_2A_2$ are $-v$ and $v$ by \eqref{eq:cones}.
In particular, both $A_j$ are traceless. In terms of the Pauli matrices
\[
 \sigma_1=\begin{pmatrix}0&1\\1&0\end{pmatrix},\quad
 \sigma_2=\begin{pmatrix}0&-\ii\\\ii&0\end{pmatrix},\quad
 \sigma_3=\begin{pmatrix}1&0\\0&-1\end{pmatrix},
\]
write $A_j=a_j\cdot\boldsymbol\sigma$, $a_j\in\R^3$.
The eigenvalue condition says
\[
 |q_1a_1+q_2a_2|^2=v^2(q_1^2+q_2^2),\qquad
 a_j\cdot a_\ell=v^2\delta_{j\ell}.
\]
A constant unitary change of basis in $\C^2$ therefore puts
\eqref{eq:linear-matrix} into the form
\begin{equation}\label{eq:dirac-matrix}
 h(q)=\mu_{\mathrm F}I+v(q_1\sigma_1+q_2\sigma_2)+O(r^2).
\end{equation}
One elementary way to see this change of basis is to diagonalize
$A_1/v$, use $A_1A_2+A_2A_1=0$ to make $A_2/v$ off diagonal, and then
adjust the phase of one basis vector. A further fixed unitary rotation
gives the displayed pair of Pauli matrices.

Write $h(q)=a_0(q)I+d(q)\cdot\boldsymbol\sigma$.
Equation \eqref{eq:dirac-matrix} implies
$d(q)=v(q_1,q_2,0)+O(r^2)$, so the two eigenprojections obey
\begin{equation}\label{eq:pi-leading}
 \pi_\pm(q)=\tfrac12(I\pm\widehat d(q)\cdot\boldsymbol\sigma)
   =\tfrac12(I\pm\omega_1\sigma_1\pm\omega_2\sigma_2)+O(r).
\end{equation}
All estimates are uniform in $\omega$, because the Taylor remainders
are uniform and $|d(q)|\geq cr$.

The compressed physical velocity, rather than the derivative of an
arbitrarily chosen eigenvector, is the relevant observable. For any
$e\in\R^2$ it satisfies
\begin{equation}\label{eq:compressed-velocity}
 V_e(q):=W(q)^*((p+k_i+q)\cdot e)W(q)
       =v(e_1\sigma_1+e_2\sigma_2)+O(r).
\end{equation}
This follows directly from the $H^2$ analyticity of $W$ and
\eqref{eq:linear-matrix}.
Using \eqref{eq:pi-leading} and \eqref{eq:compressed-velocity},
\[
 G_i^{\alpha\beta}(q)
     =\Rea\Tr_{\C^2}\bigl(\pi_-(q)V_{e_\alpha}(q)
                              \pi_+(q)V_{e_\beta}(q)\bigr).
\]
For $a,b,n\in\R^3$ with $|n|=1$, Pauli multiplication gives
\begin{equation}\label{eq:pauli-trace}
 \Rea\Tr\!\left[
     \frac{I-n\cdot\boldsymbol\sigma}{2}(a\cdot\boldsymbol\sigma)
     \frac{I+n\cdot\boldsymbol\sigma}{2}(b\cdot\boldsymbol\sigma)
             \right]
       =a\cdot b-(a\cdot n)(b\cdot n).
\end{equation}
Taking $n=(\omega_1,\omega_2,0)$ and
$a=v(e_{\alpha,1},e_{\alpha,2},0)$,
$b=v(e_{\beta,1},e_{\beta,2},0)$ proves
\eqref{eq:local-matrix}. Formula \eqref{eq:local-gap} already follows
from \eqref{eq:cones}, or from $\Delta_i(q)=2|d(q)|$.
\end{proof}

\section{The long-time Ces\`aro limit}

Fubini's theorem is applicable to \eqref{eq:response} on every finite
time interval, because its absolute value is dominated by $t$ times
the integrable sum in \eqref{eq:integrable-sum}. Therefore
\begin{equation}\label{eq:averaged-response}
 \frac1T\int_0^T\sigma_{\alpha\beta}(t)\dd t
   =\frac{2}{(2\pi)^2}\int_{\cB}
       \sum_{n\leq N<m}\Rea c_{nm}^{\alpha\beta}(k)
       \frac{1-\cos(T\Delta_{mn}(k))}
            {T\Delta_{mn}(k)^3}\dd k.
\end{equation}

Choose disjoint small disks $D_i=\{k_i+q:|q|<\delta\}$ in the
quasimomentum torus. Shrink $\delta$ so that the two-band cluster is
isolated on each disk and Lemma~\ref{lem:cone} applies. On the complement
of these disks every occupied-to-unoccupied gap is bounded below by
a positive constant. Within $D_i$, every such gap other than the one
with $(n,m)=(N,N+1)$ is also bounded below by a positive constant:
this is exactly the separation from $\lambda_{N-1}$ and
$\lambda_{N+2}$ in \eqref{eq:double}. Compactness gives a common
positive bound for all these remaining transitions. Thus
\eqref{eq:matrix-sum-bound} and $|1-\cos s|\leq2$ show that their
entire contribution to \eqref{eq:averaged-response} is $O(T^{-1})$.

It remains to evaluate, for each $i$,
\begin{equation}\label{eq:cone-integral}
 I_i(T)=\frac{2}{(2\pi)^2}\int_{|q|<\delta}
   G_i^{\alpha\beta}(q)
     \frac{1-\cos(T\Delta_i(q))}{T\Delta_i(q)^3}\dd q.
\end{equation}
Change variables $q=z/T$. The Jacobian cancels the powers of $T$:
\begin{equation}\label{eq:rescaled-integral}
 I_i(T)=\frac{2}{(2\pi)^2}\int_{\R^2}
   \one_{\{|z|<T\delta\}}G_i^{\alpha\beta}(z/T)
     \frac{1-\cos(T\Delta_i(z/T))}
          {(T\Delta_i(z/T))^3}\dd z.
\end{equation}
For each $z\ne0$, Lemma~\ref{lem:cone} gives
\[
 T\Delta_i(z/T)\longrightarrow2v_i|z|,\qquad
 G_i^{\alpha\beta}(z/T)\longrightarrow
 v_i^2\left(e_\alpha\cdot e_\beta
    -(e_\alpha\cdot\widehat z)(e_\beta\cdot\widehat z)\right).
\]
Moreover, on the support of the integrand,
$c|z|\leq T\Delta_i(z/T)\leq C|z|$ and $|G_i^{\alpha\beta}|\leq C$.
Since $1-\cos s\leq\min(s^2/2,2)$, the absolute value of the
integrand in \eqref{eq:rescaled-integral} is bounded by
\begin{equation}\label{eq:dominator}
 C\begin{cases}
        |z|^{-1},&0<|z|\leq1,\\
        |z|^{-3},&|z|>1.
   \end{cases}
\end{equation}
This is integrable on $\R^2$. Dominated convergence is therefore
applicable; in particular, it treats the nonlinear error in the gap
inside the oscillatory factor without any differentiation of that factor.

Writing $z=r\omega(\theta)$ in the limiting integral gives
\begin{align}
 \lim_{T\to\infty}I_i(T)
  =\frac{2v_i^2}{(2\pi)^2}
   &\int_0^{2\pi}\left(e_\alpha\cdot e_\beta
       -(e_\alpha\cdot\omega(\theta))
        (e_\beta\cdot\omega(\theta))\right)\dd\theta
       \notag\\[-2pt]
   &\hspace{12mm}\times
     \int_0^\infty\frac{1-\cos(2v_i r)}{(2v_i r)^3}
                          r\dd r.\label{eq:limiting-integrals}
\end{align}
The angular integral equals $\pi(e_\alpha\cdot e_\beta)$, because
$\int_0^{2\pi}\omega_j\omega_\ell\dd\theta=\pi\delta_{j\ell}$.
The radial integral, with $s=2v_i r$, equals
\[
 \frac1{4v_i^2}\int_0^\infty\frac{1-\cos s}{s^2}\dd s
       =\frac{\pi}{8v_i^2}.
\]
For a direct verification of the last equality, Tonelli's theorem and
$s^{-2}=\int_0^\infty a e^{-as}\dd a$ yield
\[
 \int_0^\infty\frac{1-\cos s}{s^2}\dd s
  =\int_0^\infty a\left(\frac1a-\frac{a}{a^2+1}\right)\dd a
  =\int_0^\infty\frac{\dd a}{1+a^2}=\frac\pi2.
\]
Consequently all dependence on the cone speed cancels, and
\begin{equation}\label{eq:one-cone}
 \lim_{T\to\infty}I_i(T)
    =\frac{2v_i^2}{(2\pi)^2}\,
        \pi(e_\alpha\cdot e_\beta)\,
        \frac{\pi}{8v_i^2}
    =\frac1{16}\,e_\alpha\cdot e_\beta.
\end{equation}
Summing \eqref{eq:one-cone} over the $M$ crossings and adding the
$O(T^{-1})$ remainder proves \eqref{eq:main}. Together with
Lemmas~\ref{lem:domains} and \ref{lem:gauge}, and
Propositions~\ref{prop:sumrule} and \ref{prop:response}, this completes
the proof of Theorem~\ref{thm:main}. \hfill$\square$

\section{Consequences and scope}

The value $1/16$ is the contribution of a single cone for one spin
species with the normalization \eqref{eq:trace-definition} and
\eqref{eq:volumes}. Thus two cones give $1/8$ for the spinless model;
two independent identical spin copies give $1/4$. Each cone is counted
once on the reciprocal torus, even if a drawing of a Brillouin cell
places representatives on several boundary edges.

Time-reversal symmetry, supplied here by realness of $V$, is used to
cancel the imaginary, antisymmetric part of the response in
\eqref{eq:full-response}. Isotropy is used in
\eqref{eq:dirac-matrix}--\eqref{eq:local-matrix}; the cone velocities
themselves need not be equal at different crossings.

The averaging in \eqref{eq:main} produces the integrable kernel
\eqref{eq:dominator}. The proof therefore requires no non-stationary
phase assumption on the other bands and makes no claim that the
unaveraged response converges at large time. Likewise, it does not
interchange the small-field and long-time limits. If the Fermi gap is
everywhere positive (the case of no crossings), the same proof gives
zero Ces\`aro limit directly from \eqref{eq:averaged-response}.

\begin{thebibliography}{9}
\bibitem{RS2}
M.~Reed and B.~Simon,
\emph{Methods of Modern Mathematical Physics, II:
Fourier Analysis, Self-Adjointness},
Academic Press, 1975.

\bibitem{RS4}
M.~Reed and B.~Simon,
\emph{Methods of Modern Mathematical Physics, IV:
Analysis of Operators},
Academic Press, 1978.
\end{thebibliography}

\end{document}
